%%%%%%%%%%%%%%%%%%%%%%%%%%%%%%%%%%%%%%%%%%%%%%%%%%%%%%%%%%%%
% 2.4 COMPUTATIONAL NUMBER THEORY
% The Composition of Advanced Mathematics
%%%%%%%%%%%%%%%%%%%%%%%%%%%%%%%%%%%%%%%%%%%%%%%%%%%%%%%%%%%%

\section{Computational Number Theory}
\label{sec:computational-number-theory}

\prerequisite{
Elementary Number Theory from \cref{sec:elementary-number-theory}, including
divisibility, greatest common divisors, prime numbers, congruences, modular
arithmetic, Euler's theorem, and the Chinese Remainder Theorem. Familiarity
with algorithms and basic programming is helpful but not required.
}

\learningobjectives{
By the end of this section, the reader should be able to:
\begin{itemize}
    \item explain the purpose of computational number theory;
    \item distinguish mathematical existence from computational feasibility;
    \item analyze arithmetic algorithms using basic complexity ideas;
    \item implement and apply the Euclidean and extended Euclidean algorithms;
    \item compute modular inverses efficiently;
    \item perform modular exponentiation using repeated squaring;
    \item solve systems of congruences computationally;
    \item generate primes using sieve algorithms;
    \item distinguish deterministic and probabilistic primality testing;
    \item apply Fermat and Miller--Rabin primality tests;
    \item explain pseudoprimes and Carmichael numbers;
    \item distinguish primality testing from integer factorization;
    \item understand trial division, Fermat factorization, Pollard's rho method,
    and the quadratic sieve conceptually;
    \item explain why the number field sieve is important for large integer
    factorization;
    \item compute multiplicative arithmetic functions from prime factorizations;
    \item understand discrete logarithm problems;
    \item apply the baby-step giant-step algorithm conceptually;
    \item understand computational aspects of quadratic residues;
    \item apply the Legendre and Jacobi symbols algorithmically;
    \item understand the computational role of continued fractions;
    \item recognize the mathematical foundations of RSA and Diffie--Hellman;
    \item distinguish computational hardness assumptions from mathematical
    impossibility;
    \item understand how computational experiments can guide number-theoretic
    conjectures without replacing proof.
\end{itemize}
}

%%%%%%%%%%%%%%%%%%%%%%%%%%%%%%%%%%%%%%%%%%%%%%%%%%%%%%%%%%%%
\subsection{What Is Computational Number Theory?}
\label{subsec:what-is-computational-number-theory}
%%%%%%%%%%%%%%%%%%%%%%%%%%%%%%%%%%%%%%%%%%%%%%%%%%%%%%%%%%%%

Number theory asks questions about integers, primes, divisibility,
congruences, and Diophantine equations.

Computational number theory asks an additional question:

\begin{center}
\emph{How can these quantities actually be computed efficiently?}
\end{center}

\begin{definitionbox}{Computational Number Theory}
\emph{Computational number theory} is the study of algorithms for solving
number-theoretic problems and of the computational resources required by
those algorithms.
\end{definitionbox}

Typical computational problems include:

\begin{itemize}
    \item computing \(\gcd(a,b)\);
    \item finding modular inverses;
    \item evaluating \(a^n\bmod m\);
    \item determining whether a large integer is prime;
    \item factoring a large composite integer;
    \item solving congruences;
    \item computing discrete logarithms;
    \item evaluating arithmetic functions;
    \item searching for solutions to Diophantine equations.
\end{itemize}

\begin{conceptbox}
A mathematical problem can be theoretically solvable while still being
computationally difficult.

Computational number theory therefore studies both

\[ \boxed{\text{correctness}} \]

and

\[ \boxed{\text{efficiency}}. \]
\end{conceptbox}

%%%%%%%%%%%%%%%%%%%%%%%%%%%%%%%%%%%%%%%%%%%%%%%%%%%%%%%%%%%%
\subsection{Algorithms}
\label{subsec:number-theoretic-algorithms}
%%%%%%%%%%%%%%%%%%%%%%%%%%%%%%%%%%%%%%%%%%%%%%%%%%%%%%%%%%%%

\begin{definitionbox}{Algorithm}
An \emph{algorithm} is a finite, precisely specified procedure for solving a
problem or computing an output from given input.
\end{definitionbox}

For example, trial division gives an algorithm for determining whether a
positive integer \(n\) has a divisor.

One could test

\[ 2,3,4,\ldots,n-1. \]

This procedure is correct, but it is unnecessarily inefficient.

A better observation is that if \(n\) is composite, then it has a divisor
not exceeding

\[ \sqrt n. \]

Thus only potential divisors through \(\sqrt n\) need to be tested.

%%%%%%%%%%%%%%%%%%%%%%%%%%%%%%%%%%%%%%%%%%%%%%%%%%%%%%%%%%%%
\subsection{Input Size}
\label{subsec:integer-input-size}
%%%%%%%%%%%%%%%%%%%%%%%%%%%%%%%%%%%%%%%%%%%%%%%%%%%%%%%%%%%%

When measuring the complexity of an integer algorithm, the size of an integer
\(n\) is not usually measured by \(n\) itself.

Instead, it is measured by the number of digits needed to represent \(n\).

In binary, this is approximately

\[ \log_2 n. \]

\begin{definitionbox}{Bit Length}
For a positive integer \(n\), its \emph{bit length} is

\[ \ell(n)=\lfloor\log_2 n\rfloor+1. \]
\end{definitionbox}

The symbol

\[ \ell \]

is lowercase script-style \emph{ell} and here denotes length.

For example,

\[ 13=(1101)_2, \]

so

\[ \ell(13)=4. \]

%%%%%%%%%%%%%%%%%%%%%%%%%%%%%%%%%%%%%%%%%%%%%%%%%%%%%%%%%%%%
\subsection{Why Bit Length Matters}
\label{subsec:why-bit-length-matters}
%%%%%%%%%%%%%%%%%%%%%%%%%%%%%%%%%%%%%%%%%%%%%%%%%%%%%%%%%%%%

Suppose an algorithm requires approximately

\[ \sqrt n \]

operations.

Since

\[ n\approx2^{\ell(n)}, \]

we have

\[ \sqrt n\approx2^{\ell(n)/2}. \]

Thus an algorithm that appears to grow like a square root of the integer may
actually grow exponentially in the number of input bits.

\begin{importantbox}
For integer algorithms, computational complexity should generally be measured
relative to the number of bits in the input, not merely the numerical value
of the integer.
\end{importantbox}

%%%%%%%%%%%%%%%%%%%%%%%%%%%%%%%%%%%%%%%%%%%%%%%%%%%%%%%%%%%%
\subsection{Basic Complexity Notation}
\label{subsec:number-theory-complexity-notation}
%%%%%%%%%%%%%%%%%%%%%%%%%%%%%%%%%%%%%%%%%%%%%%%%%%%%%%%%%%%%

Asymptotic notation provides a convenient way to compare algorithm growth.

If

\[ T(n)=O(f(n)), \]

then the running time is bounded above by a constant multiple of \(f(n)\) for
sufficiently large \(n\).

Common growth classes include

\[
O(1),\qquad O(\log n),\qquad O(n),\qquad O(n\log n),\qquad O(n^2),
\qquad O(2^n).
\]

A rough hierarchy is

\[
O(1)
\prec
O(\log n)
\prec
O(n)
\prec
O(n\log n)
\prec
O(n^2)
\prec
O(2^n).
\]

The symbol

\[ \prec \]

is read ``grows more slowly than'' in this informal comparison.

%%%%%%%%%%%%%%%%%%%%%%%%%%%%%%%%%%%%%%%%%%%%%%%%%%%%%%%%%%%%
\subsection{The Euclidean Algorithm}
\label{subsec:computational-euclidean-algorithm}
%%%%%%%%%%%%%%%%%%%%%%%%%%%%%%%%%%%%%%%%%%%%%%%%%%%%%%%%%%%%

The Euclidean algorithm is one of the oldest and most important efficient
algorithms in mathematics.

It is based on the identity

\[ \gcd(a,b)=\gcd(b,a\bmod b). \]

Repeated division produces

\[
a=bq_1+r_1,\qquad
b=r_1q_2+r_2,\qquad
r_1=r_2q_3+r_3,\qquad\ldots
\]

until a remainder becomes zero.

The last nonzero remainder is

\[ \gcd(a,b). \]

%%%%%%%%%%%%%%%%%%%%%%%%%%%%%%%%%%%%%%%%%%%%%%%%%%%%%%%%%%%%
\subsection{Euclidean Algorithm Flow}
\label{subsec:euclidean-algorithm-flow}
%%%%%%%%%%%%%%%%%%%%%%%%%%%%%%%%%%%%%%%%%%%%%%%%%%%%%%%%%%%%

\begin{centereddiagram}
\begin{tikzpicture}[
    node distance=8mm,
    box/.style={draw,rounded corners,align=center,minimum width=3.5cm,
    minimum height=0.8cm}
]
\node[box] (start) {Start with \(a>b>0\)};
\node[box,below=of start] (divide) {Compute \(a=bq+r\)};
\node[box,below=of divide] (check) {Is \(r=0\)?};
\node[box,below left=10mm and 8mm of check] (replace) {Replace\\\(a\leftarrow b,\ b\leftarrow r\)};
\node[box,below right=10mm and 8mm of check] (done) {\(\gcd(a,b)=b\)};

\draw[-{Latex[length=2.5mm]}] (start)--(divide);
\draw[-{Latex[length=2.5mm]}] (divide)--(check);
\draw[-{Latex[length=2.5mm]}] (check)--node[left]{No}(replace);
\draw[-{Latex[length=2.5mm]}] (check)--node[right]{Yes}(done);
\draw[-{Latex[length=2.5mm]}] (replace.west) -| ++(-0.8,0) |- (divide.west);
\end{tikzpicture}
\end{centereddiagram}

%%%%%%%%%%%%%%%%%%%%%%%%%%%%%%%%%%%%%%%%%%%%%%%%%%%%%%%%%%%%
\subsection{Worked Example: Euclidean Algorithm}
\label{subsec:worked-euclidean-computation}
%%%%%%%%%%%%%%%%%%%%%%%%%%%%%%%%%%%%%%%%%%%%%%%%%%%%%%%%%%%%

\begin{workedexamplebox}{Computing a Greatest Common Divisor}

Compute

\[ \gcd(252,198). \]

\begin{tabularx}{\textwidth}{@{}>{\raggedright\arraybackslash}p{0.36\textwidth}X@{}}
Divide \(252\) by \(198\). & \(252=198(1)+54\)\\
Divide \(198\) by \(54\). & \(198=54(3)+36\)\\
Divide \(54\) by \(36\). & \(54=36(1)+18\)\\
Divide \(36\) by \(18\). & \(36=18(2)+0\)
\end{tabularx}

The last nonzero remainder is \(18\).

\begin{answerbox}
\[ \boxed{\gcd(252,198)=18.} \]
\end{answerbox}
\end{workedexamplebox}

%%%%%%%%%%%%%%%%%%%%%%%%%%%%%%%%%%%%%%%%%%%%%%%%%%%%%%%%%%%%
\subsection{The Extended Euclidean Algorithm}
\label{subsec:extended-euclidean-algorithm}
%%%%%%%%%%%%%%%%%%%%%%%%%%%%%%%%%%%%%%%%%%%%%%%%%%%%%%%%%%%%

The Euclidean algorithm can be extended to produce integers \(x\) and \(y\)
satisfying Bézout's identity.

\begin{theorem}[Bézout's Identity]
For integers \(a,b\), not both zero, there exist integers \(x,y\) such that

\[ \boxed{ax+by=\gcd(a,b).} \]
\end{theorem}

The extended Euclidean algorithm computes these coefficients efficiently.

%%%%%%%%%%%%%%%%%%%%%%%%%%%%%%%%%%%%%%%%%%%%%%%%%%%%%%%%%%%%
\subsection{Worked Example: Extended Euclidean Algorithm}
\label{subsec:worked-extended-euclidean}
%%%%%%%%%%%%%%%%%%%%%%%%%%%%%%%%%%%%%%%%%%%%%%%%%%%%%%%%%%%%

Using the previous divisions,

\[ 18=54-36. \]

Since

\[ 36=198-3(54), \]

we obtain

\[ 18=54-(198-3(54))=4(54)-198. \]

Since

\[ 54=252-198, \]

we obtain

\[ 18=4(252-198)-198=4(252)-5(198). \]

Therefore

\[ \boxed{18=4(252)-5(198).} \]

Thus one Bézout pair is

\[ x=4,\qquad y=-5. \]

%%%%%%%%%%%%%%%%%%%%%%%%%%%%%%%%%%%%%%%%%%%%%%%%%%%%%%%%%%%%
\subsection{Computing Modular Inverses}
\label{subsec:computing-modular-inverses}
%%%%%%%%%%%%%%%%%%%%%%%%%%%%%%%%%%%%%%%%%%%%%%%%%%%%%%%%%%%%

An integer \(a\) has a multiplicative inverse modulo \(m\) precisely when

\[ \gcd(a,m)=1. \]

If the extended Euclidean algorithm gives

\[ ax+my=1, \]

then reducing modulo \(m\) gives

\[ ax\equiv1\pmod m. \]

Therefore

\[ \boxed{a^{-1}\equiv x\pmod m.} \]

%%%%%%%%%%%%%%%%%%%%%%%%%%%%%%%%%%%%%%%%%%%%%%%%%%%%%%%%%%%%
\subsection{Worked Example: Modular Inverse}
\label{subsec:worked-modular-inverse}
%%%%%%%%%%%%%%%%%%%%%%%%%%%%%%%%%%%%%%%%%%%%%%%%%%%%%%%%%%%%

\begin{workedexamplebox}{Finding a Modular Inverse}

Find

\[ 17^{-1}\pmod{43}. \]

\begin{tabularx}{\textwidth}{@{}>{\raggedright\arraybackslash}p{0.36\textwidth}X@{}}
Apply Euclid. & \(43=2(17)+9\)\\
Continue. & \(17=1(9)+8\)\\
Continue. & \(9=1(8)+1\)\\
Reverse the equations. & \(1=9-8\)\\
Substitute \(8=17-9\). & \(1=2(9)-17\)\\
Substitute \(9=43-2(17)\). & \(1=2(43)-5(17)\)
\end{tabularx}

Thus

\[ -5(17)\equiv1\pmod{43}. \]

Since

\[ -5\equiv38\pmod{43}, \]

we obtain

\begin{answerbox}
\[ \boxed{17^{-1}\equiv38\pmod{43}.} \]
\end{answerbox}
\end{workedexamplebox}

%%%%%%%%%%%%%%%%%%%%%%%%%%%%%%%%%%%%%%%%%%%%%%%%%%%%%%%%%%%%
\subsection{Modular Exponentiation}
\label{subsec:modular-exponentiation}
%%%%%%%%%%%%%%%%%%%%%%%%%%%%%%%%%%%%%%%%%%%%%%%%%%%%%%%%%%%%

Computing

\[ a^n\bmod m \]

by first calculating the enormous integer \(a^n\) is usually wasteful.

Instead, multiplication and reduction can be interleaved.

A fundamental method is \emph{repeated squaring}, also called
\emph{binary exponentiation}.

For example,

\[ a^8=((a^2)^2)^2. \]

More generally, write the exponent in binary.

If

\[ n=\sum_{j=0}^{k}b_j2^j,\qquad b_j\in\{0,1\}, \]

then

\[ a^n=\prod_{j:b_j=1}a^{2^j}. \]

%%%%%%%%%%%%%%%%%%%%%%%%%%%%%%%%%%%%%%%%%%%%%%%%%%%%%%%%%%%%
\subsection{Worked Example: Repeated Squaring}
\label{subsec:worked-repeated-squaring}
%%%%%%%%%%%%%%%%%%%%%%%%%%%%%%%%%%%%%%%%%%%%%%%%%%%%%%%%%%%%

\begin{workedexamplebox}{Fast Modular Exponentiation}

Compute

\[ 7^{13}\pmod{20}. \]

Since

\[ 13=8+4+1, \]

we need \(7^1,7^4,7^8\).

\begin{tabularx}{\textwidth}{@{}>{\raggedright\arraybackslash}p{0.36\textwidth}X@{}}
Start. & \(7^2=49\equiv9\pmod{20}\)\\
Square. & \(7^4\equiv9^2=81\equiv1\pmod{20}\)\\
Square again. & \(7^8\equiv1^2\equiv1\pmod{20}\)\\
Combine. & \(7^{13}=7^8\cdot7^4\cdot7\equiv1\cdot1\cdot7\)
\end{tabularx}

\begin{answerbox}
\[ \boxed{7^{13}\equiv7\pmod{20}.} \]
\end{answerbox}
\end{workedexamplebox}

%%%%%%%%%%%%%%%%%%%%%%%%%%%%%%%%%%%%%%%%%%%%%%%%%%%%%%%%%%%%
\subsection{Square-and-Multiply Algorithm}
\label{subsec:square-and-multiply}
%%%%%%%%%%%%%%%%%%%%%%%%%%%%%%%%%%%%%%%%%%%%%%%%%%%%%%%%%%%%

Repeated squaring leads to the square-and-multiply algorithm.

For an exponent with approximately

\[ \log_2 n \]

binary digits, modular exponentiation requires only on the order of

\[ O(\log n) \]

modular squaring and multiplication steps.

This is dramatically better than performing \(n-1\) repeated
multiplications.

\begin{importantbox}
Efficient modular exponentiation is one of the fundamental computational
operations underlying modern number-theoretic cryptography.
\end{importantbox}

%%%%%%%%%%%%%%%%%%%%%%%%%%%%%%%%%%%%%%%%%%%%%%%%%%%%%%%%%%%%
\subsection{Solving Linear Congruences Algorithmically}
\label{subsec:computational-linear-congruences}
%%%%%%%%%%%%%%%%%%%%%%%%%%%%%%%%%%%%%%%%%%%%%%%%%%%%%%%%%%%%

Consider

\[ ax\equiv b\pmod m. \]

Let

\[ d=\gcd(a,m). \]

A solution exists precisely when

\[ d\mid b. \]

If \(d=1\), compute

\[ a^{-1}\pmod m \]

using the extended Euclidean algorithm and obtain

\[ x\equiv a^{-1}b\pmod m. \]

Thus solving a linear congruence reduces computationally to a greatest common
divisor calculation.

%%%%%%%%%%%%%%%%%%%%%%%%%%%%%%%%%%%%%%%%%%%%%%%%%%%%%%%%%%%%
\subsection{Computational Chinese Remainder Theorem}
\label{subsec:computational-crt}
%%%%%%%%%%%%%%%%%%%%%%%%%%%%%%%%%%%%%%%%%%%%%%%%%%%%%%%%%%%%

Suppose

\[
x\equiv a_1\pmod{m_1},\qquad
x\equiv a_2\pmod{m_2},
\]

where

\[ \gcd(m_1,m_2)=1. \]

Let

\[ M=m_1m_2,\qquad M_1=\frac{M}{m_1},\qquad M_2=\frac{M}{m_2}. \]

Find

\[ y_1\equiv M_1^{-1}\pmod{m_1},\qquad y_2\equiv M_2^{-1}\pmod{m_2}. \]

Then

\[ \boxed{x\equiv a_1M_1y_1+a_2M_2y_2\pmod M.} \]

%%%%%%%%%%%%%%%%%%%%%%%%%%%%%%%%%%%%%%%%%%%%%%%%%%%%%%%%%%%%
\subsection{Worked Example: Chinese Remainder Computation}
\label{subsec:worked-computational-crt}
%%%%%%%%%%%%%%%%%%%%%%%%%%%%%%%%%%%%%%%%%%%%%%%%%%%%%%%%%%%%

\begin{workedexamplebox}{Solving a Congruence System}

Solve

\[ x\equiv2\pmod3,\qquad x\equiv3\pmod5. \]

\begin{tabularx}{\textwidth}{@{}>{\raggedright\arraybackslash}p{0.38\textwidth}X@{}}
Compute the product. & \(M=3(5)=15\)\\
Compute partial products. & \(M_1=5,\quad M_2=3\)\\
Invert \(5\) modulo \(3\). & \(5^{-1}\equiv2\pmod3\)\\
Invert \(3\) modulo \(5\). & \(3^{-1}\equiv2\pmod5\)\\
Construct the solution. & \(x\equiv2(5)(2)+3(3)(2)=38\pmod{15}\)\\
Reduce. & \(38\equiv8\pmod{15}\)
\end{tabularx}

\begin{answerbox}
\[ \boxed{x\equiv8\pmod{15}.} \]
\end{answerbox}
\end{workedexamplebox}

%%%%%%%%%%%%%%%%%%%%%%%%%%%%%%%%%%%%%%%%%%%%%%%%%%%%%%%%%%%%
\subsection{Prime Generation}
\label{subsec:prime-generation}
%%%%%%%%%%%%%%%%%%%%%%%%%%%%%%%%%%%%%%%%%%%%%%%%%%%%%%%%%%%%

A basic computational task is to generate all primes up to a given bound.

Testing each integer independently wastes information.

Sieve algorithms eliminate many composites simultaneously.

%%%%%%%%%%%%%%%%%%%%%%%%%%%%%%%%%%%%%%%%%%%%%%%%%%%%%%%%%%%%
\subsection{Sieve of Eratosthenes}
\label{subsec:computational-sieve-eratosthenes}
%%%%%%%%%%%%%%%%%%%%%%%%%%%%%%%%%%%%%%%%%%%%%%%%%%%%%%%%%%%%

To find all primes not exceeding \(N\):

\begin{enumerate}
    \item List the integers from \(2\) through \(N\).
    \item Begin with \(p=2\).
    \item Mark every multiple of \(p\) larger than \(p\) as composite.
    \item Move to the smallest unmarked integer larger than \(p\).
    \item Repeat until \(p^2>N\).
    \item The remaining unmarked integers are prime.
\end{enumerate}

%%%%%%%%%%%%%%%%%%%%%%%%%%%%%%%%%%%%%%%%%%%%%%%%%%%%%%%%%%%%
\subsection{Visual Example of the Sieve}
\label{subsec:sieve-visual}
%%%%%%%%%%%%%%%%%%%%%%%%%%%%%%%%%%%%%%%%%%%%%%%%%%%%%%%%%%%%

\begin{centereddiagram}
\begin{tikzpicture}[
    cell/.style={draw,minimum width=0.72cm,minimum height=0.55cm}
]
\foreach \n in {1,...,30}{
    \pgfmathtruncatemacro{\row}{int((\n-1)/10)}
    \pgfmathtruncatemacro{\col}{mod(\n-1,10)}
    \node[cell] (n\n) at (\col*0.72,-\row*0.55) {\scriptsize \n};
}
\foreach \p in {2,3,5,7,11,13,17,19,23,29}{
    \node[draw=MathBlue,line width=0.9pt,rounded corners,
    fit=(n\p),inner sep=0pt] {};
}
\end{tikzpicture}
\end{centereddiagram}

The highlighted numbers are the primes through \(30\).

%%%%%%%%%%%%%%%%%%%%%%%%%%%%%%%%%%%%%%%%%%%%%%%%%%%%%%%%%%%%
\subsection{Why the Sieve Stops at \(\sqrt N\)}
\label{subsec:sieve-square-root}
%%%%%%%%%%%%%%%%%%%%%%%%%%%%%%%%%%%%%%%%%%%%%%%%%%%%%%%%%%%%

If a composite integer \(n\leq N\) has a factorization

\[ n=ab, \]

then at least one factor satisfies

\[ a\leq\sqrt n\leq\sqrt N. \]

Therefore every composite number through \(N\) has already been marked once
all primes through

\[ \sqrt N \]

have been processed.

%%%%%%%%%%%%%%%%%%%%%%%%%%%%%%%%%%%%%%%%%%%%%%%%%%%%%%%%%%%%
\subsection{Primality Testing}
\label{subsec:primality-testing}
%%%%%%%%%%%%%%%%%%%%%%%%%%%%%%%%%%%%%%%%%%%%%%%%%%%%%%%%%%%%

\begin{definitionbox}{Primality Test}
A \emph{primality test} is an algorithm that determines whether a given
integer \(n>1\) is prime.
\end{definitionbox}

Primality testing and factorization are different problems.

A primality test answers:

\begin{center}
``Is \(n\) prime?''
\end{center}

A factorization algorithm asks:

\begin{center}
``If \(n\) is composite, what are its prime factors?''
\end{center}

The second problem is generally much harder computationally.

%%%%%%%%%%%%%%%%%%%%%%%%%%%%%%%%%%%%%%%%%%%%%%%%%%%%%%%%%%%%
\subsection{Trial Division}
\label{subsec:trial-division-primality}
%%%%%%%%%%%%%%%%%%%%%%%%%%%%%%%%%%%%%%%%%%%%%%%%%%%%%%%%%%%%

The simplest primality test checks whether \(n\) is divisible by any prime

\[ p\leq\sqrt n. \]

If no such prime divides \(n\), then \(n\) is prime.

This works well for small numbers but becomes inefficient for very large
integers.

%%%%%%%%%%%%%%%%%%%%%%%%%%%%%%%%%%%%%%%%%%%%%%%%%%%%%%%%%%%%
\subsection{Fermat's Primality Test}
\label{subsec:fermat-primality-test}
%%%%%%%%%%%%%%%%%%%%%%%%%%%%%%%%%%%%%%%%%%%%%%%%%%%%%%%%%%%%

Fermat's Little Theorem states that if \(p\) is prime and

\[ \gcd(a,p)=1, \]

then

\[ a^{p-1}\equiv1\pmod p. \]

This suggests a computational test.

Given an odd integer \(n\), choose a base \(a\) and compute

\[ a^{n-1}\pmod n. \]

If

\[ a^{n-1}\not\equiv1\pmod n, \]

then \(n\) is definitely composite.

%%%%%%%%%%%%%%%%%%%%%%%%%%%%%%%%%%%%%%%%%%%%%%%%%%%%%%%%%%%%
\subsection{Fermat Witnesses and Liars}
\label{subsec:fermat-witnesses}
%%%%%%%%%%%%%%%%%%%%%%%%%%%%%%%%%%%%%%%%%%%%%%%%%%%%%%%%%%%%

\begin{definitionbox}{Fermat Witness}
If

\[ a^{n-1}\not\equiv1\pmod n, \]

then \(a\) is called a \emph{Fermat witness} to the compositeness of \(n\).
\end{definitionbox}

If \(n\) is composite but nevertheless satisfies

\[ a^{n-1}\equiv1\pmod n, \]

then \(a\) is sometimes called a \emph{Fermat liar} for \(n\).

Thus passing a Fermat test does not prove primality.

%%%%%%%%%%%%%%%%%%%%%%%%%%%%%%%%%%%%%%%%%%%%%%%%%%%%%%%%%%%%
\subsection{Pseudoprimes}
\label{subsec:pseudoprimes}
%%%%%%%%%%%%%%%%%%%%%%%%%%%%%%%%%%%%%%%%%%%%%%%%%%%%%%%%%%%%

\begin{definitionbox}{Pseudoprime to Base \(a\)}
A composite integer \(n\) with

\[ \gcd(a,n)=1 \]

is a \emph{pseudoprime to base \(a\)} if

\[ a^{n-1}\equiv1\pmod n. \]
\end{definitionbox}

For example,

\[ 341=11\cdot31 \]

is composite, yet

\[ 2^{340}\equiv1\pmod{341}. \]

Thus \(341\) is a pseudoprime to base \(2\).

%%%%%%%%%%%%%%%%%%%%%%%%%%%%%%%%%%%%%%%%%%%%%%%%%%%%%%%%%%%%
\subsection{Carmichael Numbers}
\label{subsec:carmichael-numbers-computational}
%%%%%%%%%%%%%%%%%%%%%%%%%%%%%%%%%%%%%%%%%%%%%%%%%%%%%%%%%%%%

Some composite integers pass Fermat's test for every base relatively prime to
them.

\begin{definitionbox}{Carmichael Number}
A composite integer \(n\) is a \emph{Carmichael number} if

\[ a^{n-1}\equiv1\pmod n \]

for every integer \(a\) satisfying

\[ \gcd(a,n)=1. \]
\end{definitionbox}

The smallest Carmichael number is

\[ 561=3\cdot11\cdot17. \]

Carmichael numbers demonstrate why the Fermat test alone is insufficient.

%%%%%%%%%%%%%%%%%%%%%%%%%%%%%%%%%%%%%%%%%%%%%%%%%%%%%%%%%%%%
\subsection{Miller--Rabin Primality Testing}
\label{subsec:miller-rabin}
%%%%%%%%%%%%%%%%%%%%%%%%%%%%%%%%%%%%%%%%%%%%%%%%%%%%%%%%%%%%

The Miller--Rabin test strengthens Fermat's idea.

For an odd integer \(n>2\), write

\[ n-1=2^sd, \]

where \(d\) is odd.

The test examines powers beginning with

\[ a^d\pmod n \]

and repeatedly squares them.

For a prime \(n\), the sequence must satisfy particular conditions forced by
the structure of modular square roots of \(1\).

%%%%%%%%%%%%%%%%%%%%%%%%%%%%%%%%%%%%%%%%%%%%%%%%%%%%%%%%%%%%
\subsection{Strong Probable Primes}
\label{subsec:strong-probable-primes}
%%%%%%%%%%%%%%%%%%%%%%%%%%%%%%%%%%%%%%%%%%%%%%%%%%%%%%%%%%%%

For a chosen base \(a\), an odd integer \(n\) passes the Miller--Rabin test if

\[ a^d\equiv1\pmod n \]

or if for some

\[ 0\leq r<s, \]

we have

\[ a^{2^rd}\equiv-1\pmod n. \]

A composite number passing the test for a particular base is called a
\emph{strong pseudoprime} to that base.

%%%%%%%%%%%%%%%%%%%%%%%%%%%%%%%%%%%%%%%%%%%%%%%%%%%%%%%%%%%%
\subsection{Probabilistic Primality Testing}
\label{subsec:probabilistic-primality}
%%%%%%%%%%%%%%%%%%%%%%%%%%%%%%%%%%%%%%%%%%%%%%%%%%%%%%%%%%%%

For composite odd \(n\), at most one quarter of the admissible bases can be
strong liars for the Miller--Rabin test.

Thus independent random bases reduce the probability of repeated false
acceptance rapidly.

After \(k\) suitable independent rounds, a standard error bound is

\[ \boxed{\Pr(\text{false probable-prime result})\leq4^{-k}.} \]

The symbol

\[ \Pr \]

is read ``probability.''

\begin{warningbox}
A probabilistic primality test can provide extremely high confidence without,
in its probabilistic form, constituting an unconditional proof of primality.

For bounded input ranges, carefully chosen deterministic Miller--Rabin bases
can also give deterministic tests.
\end{warningbox}

%%%%%%%%%%%%%%%%%%%%%%%%%%%%%%%%%%%%%%%%%%%%%%%%%%%%%%%%%%%%
\subsection{Deterministic Primality Testing}
\label{subsec:deterministic-primality}
%%%%%%%%%%%%%%%%%%%%%%%%%%%%%%%%%%%%%%%%%%%%%%%%%%%%%%%%%%%%

A deterministic primality test always produces a mathematically guaranteed
answer.

One major theoretical result is the AKS primality test, which established
that primality testing belongs to deterministic polynomial time.

Its importance is largely theoretical: practical primality testing often
uses other algorithms that are considerably faster in common applications.

%%%%%%%%%%%%%%%%%%%%%%%%%%%%%%%%%%%%%%%%%%%%%%%%%%%%%%%%%%%%
\subsection{Primality Certificates}
\label{subsec:primality-certificates}
%%%%%%%%%%%%%%%%%%%%%%%%%%%%%%%%%%%%%%%%%%%%%%%%%%%%%%%%%%%%

Another strategy is to produce a certificate that can be checked efficiently.

Instead of merely reporting

\begin{center}
``\(n\) is prime,''
\end{center}

an algorithm provides auxiliary information from which primality follows.

Certificate-based methods separate:

\[
\boxed{
\text{finding a proof}
\qquad\text{from}\qquad
\text{verifying a proof}.
}
\]

This distinction appears throughout computational mathematics.

%%%%%%%%%%%%%%%%%%%%%%%%%%%%%%%%%%%%%%%%%%%%%%%%%%%%%%%%%%%%
\subsection{Integer Factorization}
\label{subsec:integer-factorization-computational}
%%%%%%%%%%%%%%%%%%%%%%%%%%%%%%%%%%%%%%%%%%%%%%%%%%%%%%%%%%%%

\begin{definitionbox}{Integer Factorization Problem}
Given a composite positive integer \(n\), the \emph{integer factorization
problem} is to determine prime numbers \(p_1,\ldots,p_r\) and positive
integers \(e_1,\ldots,e_r\) such that

\[ n=p_1^{e_1}\cdots p_r^{e_r}. \]
\end{definitionbox}

The Fundamental Theorem of Arithmetic guarantees that this factorization
exists and is unique.

It does not tell us how efficiently it can be found.

%%%%%%%%%%%%%%%%%%%%%%%%%%%%%%%%%%%%%%%%%%%%%%%%%%%%%%%%%%%%
\subsection{Trial Division Factorization}
\label{subsec:trial-division-factorization}
%%%%%%%%%%%%%%%%%%%%%%%%%%%%%%%%%%%%%%%%%%%%%%%%%%%%%%%%%%%%

The simplest factorization algorithm repeatedly tests possible prime divisors.

For each prime

\[ p\leq\sqrt n, \]

check whether

\[ p\mid n. \]

Whenever a factor is found, divide it out and continue.

This is practical for small integers but unsuitable for very large
semiprimes.

%%%%%%%%%%%%%%%%%%%%%%%%%%%%%%%%%%%%%%%%%%%%%%%%%%%%%%%%%%%%
\subsection{Fermat Factorization}
\label{subsec:fermat-factorization}
%%%%%%%%%%%%%%%%%%%%%%%%%%%%%%%%%%%%%%%%%%%%%%%%%%%%%%%%%%%%

Suppose an odd integer \(n\) can be written as

\[ n=ab. \]

If the factors are near one another, write

\[ a=x-y,\qquad b=x+y. \]

Then

\[ n=x^2-y^2. \]

Thus

\[ x^2-n=y^2. \]

Fermat factorization searches for an integer \(x\geq\sqrt n\) such that

\[ x^2-n \]

is a perfect square.

%%%%%%%%%%%%%%%%%%%%%%%%%%%%%%%%%%%%%%%%%%%%%%%%%%%%%%%%%%%%
\subsection{Worked Example: Fermat Factorization}
\label{subsec:worked-fermat-factorization}
%%%%%%%%%%%%%%%%%%%%%%%%%%%%%%%%%%%%%%%%%%%%%%%%%%%%%%%%%%%%

\begin{workedexamplebox}{Factoring by a Difference of Squares}

Factor

\[ 5959. \]

Start near

\[ \sqrt{5959}\approx77.2. \]

Try \(x=78\):

\[ 78^2-5959=6084-5959=125, \]

not a square.

Try \(x=79\):

\[ 79^2-5959=6241-5959=282. \]

Continue until \(x=80\):

\[ 80^2-5959=6400-5959=441=21^2. \]

Therefore

\[ 5959=80^2-21^2=(80-21)(80+21). \]

\begin{answerbox}
\[ \boxed{5959=59\cdot101.} \]
\end{answerbox}
\end{workedexamplebox}

%%%%%%%%%%%%%%%%%%%%%%%%%%%%%%%%%%%%%%%%%%%%%%%%%%%%%%%%%%%%
\subsection{Pollard's Rho Method}
\label{subsec:pollard-rho}
%%%%%%%%%%%%%%%%%%%%%%%%%%%%%%%%%%%%%%%%%%%%%%%%%%%%%%%%%%%%

Pollard's rho method searches for a nontrivial factor without testing all
possible divisors.

It constructs a pseudorandom sequence modulo \(n\), often using

\[ x_{k+1}\equiv x_k^2+c\pmod n. \]

If two sequence values become equal modulo an unknown factor \(p\), then

\[ p\mid x_i-x_j. \]

Consequently,

\[ \gcd(|x_i-x_j|,n) \]

may reveal a nontrivial factor of \(n\).

%%%%%%%%%%%%%%%%%%%%%%%%%%%%%%%%%%%%%%%%%%%%%%%%%%%%%%%%%%%%
\subsection{Why the Name Rho?}
\label{subsec:why-pollard-rho}
%%%%%%%%%%%%%%%%%%%%%%%%%%%%%%%%%%%%%%%%%%%%%%%%%%%%%%%%%%%%

A sequence in a finite set must eventually repeat.

Its structure resembles a tail followed by a cycle:

\begin{centereddiagram}
\begin{tikzpicture}[
    every node/.style={circle,draw,minimum size=6mm},
    >=Latex
]
\node (a) at (0,0) {};
\node (b) at (1,0) {};
\node (c) at (2,0) {};
\node (d) at (3,0) {};
\node (e) at (4,0.8) {};
\node (f) at (5,0) {};
\node (g) at (4,-0.8) {};

\draw[->] (a)--(b);
\draw[->] (b)--(c);
\draw[->] (c)--(d);
\draw[->] (d)--(e);
\draw[->] (e)--(f);
\draw[->] (f)--(g);
\draw[->] (g)--(d);
\end{tikzpicture}
\end{centereddiagram}

This tail-and-cycle shape resembles the Greek letter

\[ \rho, \]

called \emph{rho}, pronounced ``row.''

%%%%%%%%%%%%%%%%%%%%%%%%%%%%%%%%%%%%%%%%%%%%%%%%%%%%%%%%%%%%
\subsection{Cycle Detection}
\label{subsec:pollard-cycle-detection}
%%%%%%%%%%%%%%%%%%%%%%%%%%%%%%%%%%%%%%%%%%%%%%%%%%%%%%%%%%%%

Pollard's rho algorithm commonly uses Floyd's cycle-detection method.

One sequence value advances one step at a time:

\[ x\leftarrow f(x), \]

while another advances two:

\[ y\leftarrow f(f(y)). \]

At each stage compute

\[ d=\gcd(|x-y|,n). \]

If

\[ 1<d<n, \]

then \(d\) is a nontrivial factor.

%%%%%%%%%%%%%%%%%%%%%%%%%%%%%%%%%%%%%%%%%%%%%%%%%%%%%%%%%%%%
\subsection{Quadratic Sieve}
\label{subsec:quadratic-sieve}
%%%%%%%%%%%%%%%%%%%%%%%%%%%%%%%%%%%%%%%%%%%%%%%%%%%%%%%%%%%%

For larger integers, more sophisticated factorization methods search for
congruences of squares.

If

\[ x^2\equiv y^2\pmod n, \]

then

\[ n\mid(x-y)(x+y). \]

If

\[ x\not\equiv\pm y\pmod n, \]

then

\[ \gcd(x-y,n) \]

often gives a nontrivial factor.

The quadratic sieve constructs many relations whose values factor over a
small collection of primes called a \emph{factor base}.

%%%%%%%%%%%%%%%%%%%%%%%%%%%%%%%%%%%%%%%%%%%%%%%%%%%%%%%%%%%%
\subsection{Smooth Numbers}
\label{subsec:smooth-numbers}
%%%%%%%%%%%%%%%%%%%%%%%%%%%%%%%%%%%%%%%%%%%%%%%%%%%%%%%%%%%%

\begin{definitionbox}{Smooth Number}
A positive integer is called \(B\)-\emph{smooth} if all of its prime factors
are at most \(B\).
\end{definitionbox}

For example,

\[ 72=2^3\cdot3^2 \]

is \(3\)-smooth.

But

\[ 77=7\cdot11 \]

is not \(7\)-smooth because of the factor \(11\).

Smooth numbers are fundamental in modern factorization algorithms.

%%%%%%%%%%%%%%%%%%%%%%%%%%%%%%%%%%%%%%%%%%%%%%%%%%%%%%%%%%%%
\subsection{Linear Algebra in Factorization}
\label{subsec:factorization-linear-algebra}
%%%%%%%%%%%%%%%%%%%%%%%%%%%%%%%%%%%%%%%%%%%%%%%%%%%%%%%%%%%%

Suppose several smooth values have prime factorizations

\[ Q_i=\prod_{j=1}^{r}p_j^{e_{ij}}. \]

To make a product

\[ Q_{i_1}\cdots Q_{i_k} \]

a perfect square, every total exponent must be even.

Thus we consider exponent vectors modulo \(2\):

\[ (e_{i1},\ldots,e_{ir})\pmod2. \]

Finding a subset whose sum is zero becomes a linear algebra problem over

\[ \mathbf F_2. \]

\begin{conceptbox}
Modern integer factorization combines several branches of mathematics:

\[
\boxed{
\text{number theory}
+
\text{probability}
+
\text{linear algebra}
+
\text{algorithms}.
}
\]
\end{conceptbox}

%%%%%%%%%%%%%%%%%%%%%%%%%%%%%%%%%%%%%%%%%%%%%%%%%%%%%%%%%%%%
\subsection{The Number Field Sieve}
\label{subsec:number-field-sieve}
%%%%%%%%%%%%%%%%%%%%%%%%%%%%%%%%%%%%%%%%%%%%%%%%%%%%%%%%%%%%

The general number field sieve is one of the most important classical
algorithms for factoring very large integers.

Its basic philosophy is to construct smooth-number relations simultaneously
in two algebraic settings and combine them to obtain a congruence of squares
modulo the integer being factored.

The full algorithm uses ideas from:

\begin{itemize}
    \item algebraic number theory;
    \item polynomial arithmetic;
    \item smooth numbers;
    \item finite fields;
    \item sparse linear algebra.
\end{itemize}

The detailed algorithm belongs at a more advanced computational level, but
its existence illustrates the practical importance of number fields.

%%%%%%%%%%%%%%%%%%%%%%%%%%%%%%%%%%%%%%%%%%%%%%%%%%%%%%%%%%%%
\subsection{Primality Testing Versus Factorization}
\label{subsec:primality-versus-factorization}
%%%%%%%%%%%%%%%%%%%%%%%%%%%%%%%%%%%%%%%%%%%%%%%%%%%%%%%%%%%%

It is important not to confuse the two problems.

\[
\boxed{
\text{Primality testing: Is }n\text{ prime?}
}
\]

\[
\boxed{
\text{Factorization: What are the prime factors of }n\text{?}
}
\]

Efficient algorithms are known for determining primality.

No classical polynomial-time algorithm is currently known for general
integer factorization.

This distinction is fundamental in computational number theory.

%%%%%%%%%%%%%%%%%%%%%%%%%%%%%%%%%%%%%%%%%%%%%%%%%%%%%%%%%%%%
\subsection{Computing Arithmetic Functions}
\label{subsec:computing-arithmetic-functions}
%%%%%%%%%%%%%%%%%%%%%%%%%%%%%%%%%%%%%%%%%%%%%%%%%%%%%%%%%%%%

Once

\[ n=p_1^{e_1}\cdots p_r^{e_r} \]

is known, many arithmetic functions can be computed efficiently.

For example,

\[ \boxed{\tau(n)=\prod_{i=1}^{r}(e_i+1)} \]

and

\[ \boxed{\sigma(n)=\prod_{i=1}^{r}\frac{p_i^{e_i+1}-1}{p_i-1}}. \]

Euler's phi function is

\[ \boxed{\varphi(n)=n\prod_{p\mid n}\left(1-\frac1p\right).} \]

The Möbius function is

\[
\mu(n)=
\begin{cases}
0, & p^2\mid n\text{ for some prime }p,\\
(-1)^r, & n\text{ is a product of }r\text{ distinct primes}.
\end{cases}
\]

%%%%%%%%%%%%%%%%%%%%%%%%%%%%%%%%%%%%%%%%%%%%%%%%%%%%%%%%%%%%
\subsection{Computational Dependence on Factorization}
\label{subsec:functions-depend-factorization}
%%%%%%%%%%%%%%%%%%%%%%%%%%%%%%%%%%%%%%%%%%%%%%%%%%%%%%%%%%%%

These formulas reveal an important computational issue.

The formulas themselves are easy to evaluate once the prime factorization is
known.

But finding that factorization may be difficult.

Thus the computational cost of evaluating an arithmetic function can depend
strongly on what information about the input is already available.

%%%%%%%%%%%%%%%%%%%%%%%%%%%%%%%%%%%%%%%%%%%%%%%%%%%%%%%%%%%%
\subsection{Quadratic Residues}
\label{subsec:computational-quadratic-residues}
%%%%%%%%%%%%%%%%%%%%%%%%%%%%%%%%%%%%%%%%%%%%%%%%%%%%%%%%%%%%

For an odd prime \(p\), an integer \(a\) is a quadratic residue modulo \(p\)
if

\[ x^2\equiv a\pmod p \]

has a solution.

Rather than searching through all possible \(x\), Euler's criterion gives

\[ \boxed{a^{(p-1)/2}\equiv\left(\frac ap\right)\pmod p.} \]

The notation

\[ \left(\frac ap\right) \]

is the Legendre symbol, introduced in Elementary Number Theory.

It takes values

\[ -1,\quad0,\quad1. \]

%%%%%%%%%%%%%%%%%%%%%%%%%%%%%%%%%%%%%%%%%%%%%%%%%%%%%%%%%%%%
\subsection{Computing Legendre Symbols}
\label{subsec:computing-legendre-symbols}
%%%%%%%%%%%%%%%%%%%%%%%%%%%%%%%%%%%%%%%%%%%%%%%%%%%%%%%%%%%%

Quadratic reciprocity allows Legendre symbols to be evaluated without
computing large powers.

For distinct odd primes \(p,q\),

\[
\left(\frac pq\right)
\left(\frac qp\right)
=
(-1)^{\frac{(p-1)(q-1)}4}.
\]

Together with

\[ \left(\frac{-1}{p}\right)=(-1)^{(p-1)/2} \]

and

\[ \left(\frac2p\right)=(-1)^{(p^2-1)/8}, \]

this produces an efficient recursive procedure.

%%%%%%%%%%%%%%%%%%%%%%%%%%%%%%%%%%%%%%%%%%%%%%%%%%%%%%%%%%%%
\subsection{Jacobi Symbol}
\label{subsec:jacobi-symbol-computation}
%%%%%%%%%%%%%%%%%%%%%%%%%%%%%%%%%%%%%%%%%%%%%%%%%%%%%%%%%%%%

The Legendre symbol extends to odd composite denominators.

If

\[ n=p_1^{e_1}\cdots p_r^{e_r} \]

is odd, define

\[
\boxed{
\left(\frac an\right)
=
\prod_{i=1}^{r}
\left(\frac{a}{p_i}\right)^{e_i}.
}
\]

This is called the \emph{Jacobi symbol}.

A Jacobi symbol can be computed efficiently without first factoring \(n\),
using reciprocity identities and a Euclidean-algorithm-style procedure.

\begin{warningbox}
For composite \(n\),

\[ \left(\frac an\right)=1 \]

does not necessarily imply that \(a\) is a quadratic residue modulo \(n\).
\end{warningbox}

%%%%%%%%%%%%%%%%%%%%%%%%%%%%%%%%%%%%%%%%%%%%%%%%%%%%%%%%%%%%
\subsection{Modular Square Roots}
\label{subsec:modular-square-roots}
%%%%%%%%%%%%%%%%%%%%%%%%%%%%%%%%%%%%%%%%%%%%%%%%%%%%%%%%%%%%

Given a prime \(p\) and a quadratic residue \(a\), one may wish to solve

\[ x^2\equiv a\pmod p. \]

For primes satisfying

\[ p\equiv3\pmod4, \]

a particularly simple formula is

\[ \boxed{x\equiv a^{(p+1)/4}\pmod p.} \]

For general odd primes, the Tonelli--Shanks algorithm provides an efficient
method for computing modular square roots.

%%%%%%%%%%%%%%%%%%%%%%%%%%%%%%%%%%%%%%%%%%%%%%%%%%%%%%%%%%%%
\subsection{Discrete Logarithms}
\label{subsec:computational-discrete-logarithms}
%%%%%%%%%%%%%%%%%%%%%%%%%%%%%%%%%%%%%%%%%%%%%%%%%%%%%%%%%%%%

Ordinary logarithms invert exponentiation over the real numbers.

A similar problem arises in finite groups.

\begin{definitionbox}{Discrete Logarithm}
Let \(G=\langle g\rangle\) be a finite cyclic group. Given \(h\in G\), the
\emph{discrete logarithm problem} is to find an integer \(x\) satisfying

\[ \boxed{g^x=h.} \]
\end{definitionbox}

For modular arithmetic, this often takes the form

\[ g^x\equiv h\pmod p. \]

The unknown \(x\) is called a discrete logarithm of \(h\) to the base \(g\).

%%%%%%%%%%%%%%%%%%%%%%%%%%%%%%%%%%%%%%%%%%%%%%%%%%%%%%%%%%%%
\subsection{Brute-Force Discrete Logarithms}
\label{subsec:brute-force-discrete-log}
%%%%%%%%%%%%%%%%%%%%%%%%%%%%%%%%%%%%%%%%%%%%%%%%%%%%%%%%%%%%

The simplest method computes

\[ g^0,g^1,g^2,g^3,\ldots \]

until \(h\) appears.

If the group has order \(N\), this can require approximately

\[ O(N) \]

group operations.

A meet-in-the-middle strategy improves this substantially.

%%%%%%%%%%%%%%%%%%%%%%%%%%%%%%%%%%%%%%%%%%%%%%%%%%%%%%%%%%%%
\subsection{Baby-Step Giant-Step}
\label{subsec:baby-step-giant-step}
%%%%%%%%%%%%%%%%%%%%%%%%%%%%%%%%%%%%%%%%%%%%%%%%%%%%%%%%%%%%

Suppose

\[ g^x=h \]

in a cyclic group of order approximately \(N\).

Choose

\[ m=\lceil\sqrt N\rceil \]

and write

\[ x=im+j, \]

where

\[ 0\leq i,j<m. \]

Then

\[ g^{im+j}=h \]

implies

\[ g^j=h(g^{-m})^i. \]

The algorithm stores the values

\[ g^j \]

as \emph{baby steps} and compares them with the \emph{giant steps}

\[ h(g^{-m})^i. \]

The time and memory requirements are approximately

\[ O(\sqrt N). \]

%%%%%%%%%%%%%%%%%%%%%%%%%%%%%%%%%%%%%%%%%%%%%%%%%%%%%%%%%%%%
\subsection{Pollard's Rho for Discrete Logarithms}
\label{subsec:pollard-rho-discrete-log}
%%%%%%%%%%%%%%%%%%%%%%%%%%%%%%%%%%%%%%%%%%%%%%%%%%%%%%%%%%%%

A variant of Pollard's rho method can solve discrete logarithms using
approximately

\[ O(\sqrt N) \]

group operations while requiring much less memory than baby-step
giant-step.

This demonstrates an important computational tradeoff:

\[
\boxed{
\text{time}
\qquad\text{versus}\qquad
\text{memory}.
}
\]

%%%%%%%%%%%%%%%%%%%%%%%%%%%%%%%%%%%%%%%%%%%%%%%%%%%%%%%%%%%%
\subsection{Continued Fractions}
\label{subsec:computational-continued-fractions}
%%%%%%%%%%%%%%%%%%%%%%%%%%%%%%%%%%%%%%%%%%%%%%%%%%%%%%%%%%%%

Continued fractions provide highly efficient rational approximations to real
numbers.

A finite continued fraction has the form

\[
[a_0;a_1,\ldots,a_n]
=
a_0+
\cfrac{1}{
a_1+\cfrac{1}{
a_2+\cfrac{1}{
\ddots+\cfrac1{a_n}}}}.
\]

The semicolon separates the integer part \(a_0\) from the remaining partial
quotients.

%%%%%%%%%%%%%%%%%%%%%%%%%%%%%%%%%%%%%%%%%%%%%%%%%%%%%%%%%%%%
\subsection{Continued Fractions and the Euclidean Algorithm}
\label{subsec:continued-fractions-euclid}
%%%%%%%%%%%%%%%%%%%%%%%%%%%%%%%%%%%%%%%%%%%%%%%%%%%%%%%%%%%%

The Euclidean algorithm naturally produces continued fractions.

For example,

\[ \frac{43}{17}=2+\frac9{17}. \]

Continuing the Euclidean divisions gives

\[ \frac{43}{17}=[2;1,1,8]. \]

Thus the Euclidean algorithm simultaneously computes greatest common divisors
and continued-fraction expansions of rational numbers.

%%%%%%%%%%%%%%%%%%%%%%%%%%%%%%%%%%%%%%%%%%%%%%%%%%%%%%%%%%%%
\subsection{Convergents}
\label{subsec:continued-fraction-convergents}
%%%%%%%%%%%%%%%%%%%%%%%%%%%%%%%%%%%%%%%%%%%%%%%%%%%%%%%%%%%%

Truncating a continued fraction gives rational approximations called
\emph{convergents}.

If

\[ \alpha=[a_0;a_1,a_2,\ldots], \]

then its convergents are

\[
\frac{p_0}{q_0},
\frac{p_1}{q_1},
\frac{p_2}{q_2},
\ldots.
\]

Continued-fraction convergents provide exceptionally good rational
approximations and are central to computational Diophantine approximation.

%%%%%%%%%%%%%%%%%%%%%%%%%%%%%%%%%%%%%%%%%%%%%%%%%%%%%%%%%%%%
\subsection{Computational Pell Equations}
\label{subsec:computational-pell-equations}
%%%%%%%%%%%%%%%%%%%%%%%%%%%%%%%%%%%%%%%%%%%%%%%%%%%%%%%%%%%%

Continued fractions provide an effective method for solving Pell equations

\[ x^2-dy^2=1. \]

The continued-fraction expansion of

\[ \sqrt d \]

is periodic when \(d\) is a nonsquare positive integer.

Appropriate convergents yield solutions to the Pell equation.

Thus an apparently nonlinear Diophantine equation can be solved through a
systematic algorithm.

%%%%%%%%%%%%%%%%%%%%%%%%%%%%%%%%%%%%%%%%%%%%%%%%%%%%%%%%%%%%
\subsection{RSA Arithmetic}
\label{subsec:rsa-arithmetic}
%%%%%%%%%%%%%%%%%%%%%%%%%%%%%%%%%%%%%%%%%%%%%%%%%%%%%%%%%%%%

RSA is a classical public-key cryptographic construction based on modular
arithmetic.

Choose distinct large primes

\[ p,\qquad q \]

and form

\[ n=pq. \]

Then

\[ \varphi(n)=(p-1)(q-1). \]

Choose an integer \(e\) satisfying

\[ \gcd(e,\varphi(n))=1. \]

Compute \(d\) such that

\[ ed\equiv1\pmod{\varphi(n)}. \]

The pair

\[ (n,e) \]

forms the public arithmetic information, while \(d\) is derived using
information about the factorization of \(n\).

%%%%%%%%%%%%%%%%%%%%%%%%%%%%%%%%%%%%%%%%%%%%%%%%%%%%%%%%%%%%
\subsection{RSA Transformation}
\label{subsec:rsa-transformation}
%%%%%%%%%%%%%%%%%%%%%%%%%%%%%%%%%%%%%%%%%%%%%%%%%%%%%%%%%%%%

In the basic mathematical model, a message representative \(m\) is
transformed by

\[ c\equiv m^e\pmod n. \]

The inverse transformation uses

\[ m\equiv c^d\pmod n. \]

The correctness follows from modular arithmetic, ultimately using Euler's
theorem or the Chinese Remainder Theorem.

\begin{warningbox}
Textbook RSA is a mathematical model, not a secure implementation protocol.

Real cryptographic systems require carefully designed encoding, padding,
key-generation, side-channel protection, and standardized implementation
procedures.
\end{warningbox}

%%%%%%%%%%%%%%%%%%%%%%%%%%%%%%%%%%%%%%%%%%%%%%%%%%%%%%%%%%%%
\subsection{Why RSA Is Computational}
\label{subsec:why-rsa-computational}
%%%%%%%%%%%%%%%%%%%%%%%%%%%%%%%%%%%%%%%%%%%%%%%%%%%%%%%%%%%%

Several algorithms interact:

\begin{itemize}
    \item large-prime generation;
    \item probabilistic primality testing;
    \item greatest common divisor computation;
    \item modular inversion;
    \item modular exponentiation;
    \item Chinese remainder computations.
\end{itemize}

The underlying asymmetry is that multiplication

\[ p,q\longmapsto pq \]

is computationally easy, while recovering \(p\) and \(q\) from a suitably
chosen large product is believed to be computationally difficult for
classical algorithms.

%%%%%%%%%%%%%%%%%%%%%%%%%%%%%%%%%%%%%%%%%%%%%%%%%%%%%%%%%%%%
\subsection{Diffie--Hellman and Discrete Logarithms}
\label{subsec:diffie-hellman-number-theory}
%%%%%%%%%%%%%%%%%%%%%%%%%%%%%%%%%%%%%%%%%%%%%%%%%%%%%%%%%%%%

Another classical cryptographic construction uses exponentiation in a finite
cyclic group.

Given a generator \(g\), users may compute powers such as

\[ g^a,\qquad g^b,\qquad g^{ab}. \]

Exponentiation can be performed efficiently using repeated squaring.

Recovering \(a\) from

\[ g^a \]

is a discrete logarithm problem.

This creates another computational asymmetry:

\[
\boxed{
\text{modular exponentiation is easy}
\qquad
\text{while discrete logarithms may be hard}.
}
\]

%%%%%%%%%%%%%%%%%%%%%%%%%%%%%%%%%%%%%%%%%%%%%%%%%%%%%%%%%%%%
\subsection{Hardness Is Not Impossibility}
\label{subsec:hardness-not-impossibility}
%%%%%%%%%%%%%%%%%%%%%%%%%%%%%%%%%%%%%%%%%%%%%%%%%%%%%%%%%%%%

A computational problem being \emph{hard} does not mean that it is
mathematically impossible.

For example, a large composite integer certainly has a prime factorization.

The difficulty lies in finding that factorization with feasible resources.

Likewise, a discrete logarithm may exist uniquely within a cyclic group, yet
computing it may require substantial work.

%%%%%%%%%%%%%%%%%%%%%%%%%%%%%%%%%%%%%%%%%%%%%%%%%%%%%%%%%%%%
\subsection{Deterministic and Randomized Algorithms}
\label{subsec:deterministic-randomized-number-theory}
%%%%%%%%%%%%%%%%%%%%%%%%%%%%%%%%%%%%%%%%%%%%%%%%%%%%%%%%%%%%

\begin{definitionbox}{Deterministic Algorithm}
A \emph{deterministic algorithm} follows a completely determined sequence of
steps for a given input.
\end{definitionbox}

\begin{definitionbox}{Randomized Algorithm}
A \emph{randomized algorithm} uses random choices as part of its computation.
\end{definitionbox}

Randomization can improve expected running time or simplify algorithms.

Examples in computational number theory include:

\begin{itemize}
    \item Miller--Rabin primality testing;
    \item Pollard's rho factorization;
    \item random searches for suitable primes.
\end{itemize}

%%%%%%%%%%%%%%%%%%%%%%%%%%%%%%%%%%%%%%%%%%%%%%%%%%%%%%%%%%%%
\subsection{Monte Carlo and Las Vegas Algorithms}
\label{subsec:monte-carlo-las-vegas}
%%%%%%%%%%%%%%%%%%%%%%%%%%%%%%%%%%%%%%%%%%%%%%%%%%%%%%%%%%%%

Randomized algorithms are often divided into two broad types.

\begin{definitionbox}{Monte Carlo Algorithm}
A \emph{Monte Carlo algorithm} has bounded running time or predictable
performance but may have a small probability of returning an incorrect
answer.
\end{definitionbox}

\begin{definitionbox}{Las Vegas Algorithm}
A \emph{Las Vegas algorithm} always returns a correct answer, but its running
time depends on random choices.
\end{definitionbox}

Probabilistic Miller--Rabin testing is commonly viewed in a Monte Carlo
framework, while randomized factor searches can often be interpreted in a
Las Vegas framework when results are verified before acceptance.

%%%%%%%%%%%%%%%%%%%%%%%%%%%%%%%%%%%%%%%%%%%%%%%%%%%%%%%%%%%%
\subsection{Verification}
\label{subsec:number-theoretic-verification}
%%%%%%%%%%%%%%%%%%%%%%%%%%%%%%%%%%%%%%%%%%%%%%%%%%%%%%%%%%%%

Number theory frequently allows computational outputs to be verified more
easily than they were discovered.

If an algorithm claims

\[ n=ab, \]

verification requires only multiplication.

If it claims that

\[ ax+by=\gcd(a,b), \]

the identity can be checked directly.

If it proposes a solution to

\[ x^2-dy^2=1, \]

substitution verifies the result.

\begin{conceptbox}
A recurring computational principle is

\[
\boxed{
\text{finding a solution may be difficult, while checking it may be easy}.
}
\]
\end{conceptbox}

%%%%%%%%%%%%%%%%%%%%%%%%%%%%%%%%%%%%%%%%%%%%%%%%%%%%%%%%%%%%
\subsection{Computational Experiments}
\label{subsec:number-theory-computational-experiments}
%%%%%%%%%%%%%%%%%%%%%%%%%%%%%%%%%%%%%%%%%%%%%%%%%%%%%%%%%%%%

Computers allow enormous quantities of arithmetic data to be explored.

For example, one may investigate:

\begin{itemize}
    \item prime gaps;
    \item twin primes;
    \item values of arithmetic functions;
    \item pseudoprimes;
    \item quadratic residues;
    \item Diophantine solutions;
    \item class numbers;
    \item zeros of zeta and \(L\)-functions.
\end{itemize}

Patterns found computationally can suggest conjectures.

However, experimental evidence is not a proof.

%%%%%%%%%%%%%%%%%%%%%%%%%%%%%%%%%%%%%%%%%%%%%%%%%%%%%%%%%%%%
\subsection{Experiment Versus Proof}
\label{subsec:experiment-versus-proof}
%%%%%%%%%%%%%%%%%%%%%%%%%%%%%%%%%%%%%%%%%%%%%%%%%%%%%%%%%%%%

Suppose a property has been verified for the first

\[ 10^{12} \]

positive integers.

This is powerful evidence.

But it does not prove that the property holds for every positive integer.

A universal theorem requires a mathematical argument covering infinitely many
cases.

\begin{warningbox}
No finite computation can by itself prove an unrestricted statement about
all integers unless additional mathematics reduces the infinite problem to a
finite verification.
\end{warningbox}

%%%%%%%%%%%%%%%%%%%%%%%%%%%%%%%%%%%%%%%%%%%%%%%%%%%%%%%%%%%%
\subsection{Exact Integer Arithmetic}
\label{subsec:exact-integer-arithmetic}
%%%%%%%%%%%%%%%%%%%%%%%%%%%%%%%%%%%%%%%%%%%%%%%%%%%%%%%%%%%%

Number-theoretic computation usually requires exact arithmetic.

Floating-point approximations can introduce errors into calculations involving

\[ \gcd,\qquad\text{divisibility},\qquad\text{congruences}, \]

and exact factorization.

Computer algebra systems and arbitrary-precision integer libraries therefore
represent integers exactly, even when they contain thousands or millions of
digits.

%%%%%%%%%%%%%%%%%%%%%%%%%%%%%%%%%%%%%%%%%%%%%%%%%%%%%%%%%%%%
\subsection{Large Integer Multiplication}
\label{subsec:large-integer-multiplication}
%%%%%%%%%%%%%%%%%%%%%%%%%%%%%%%%%%%%%%%%%%%%%%%%%%%%%%%%%%%%

Many number-theoretic algorithms repeatedly multiply very large integers.

The elementary school multiplication algorithm for \(n\)-digit numbers uses
approximately

\[ O(n^2) \]

single-digit operations.

More advanced methods reduce this cost.

Examples include:

\begin{itemize}
    \item Karatsuba multiplication;
    \item Toom--Cook multiplication;
    \item Fourier-transform-based multiplication.
\end{itemize}

Thus even basic arithmetic becomes an algorithmic subject when integers are
sufficiently large.

%%%%%%%%%%%%%%%%%%%%%%%%%%%%%%%%%%%%%%%%%%%%%%%%%%%%%%%%%%%%
\subsection{Fast GCD and Modular Arithmetic}
\label{subsec:fast-gcd-modular}
%%%%%%%%%%%%%%%%%%%%%%%%%%%%%%%%%%%%%%%%%%%%%%%%%%%%%%%%%%%%

Efficient implementations combine several layers of algorithms.

For example, modular exponentiation requires repeated modular multiplication,
which itself requires large-integer multiplication and reduction.

Thus the true cost of

\[ a^n\bmod m \]

depends not only on the number of multiplication steps but also on the size
of the integers being multiplied.

This motivates \emph{bit complexity}, which counts operations at the level
of binary representation.

%%%%%%%%%%%%%%%%%%%%%%%%%%%%%%%%%%%%%%%%%%%%%%%%%%%%%%%%%%%%
\subsection{Precomputation and Memory}
\label{subsec:precomputation-memory}
%%%%%%%%%%%%%%%%%%%%%%%%%%%%%%%%%%%%%%%%%%%%%%%%%%%%%%%%%%%%

Algorithms can sometimes become faster by storing previously computed
information.

Examples include:

\begin{itemize}
    \item tables of small primes;
    \item precomputed modular powers;
    \item baby-step tables for discrete logarithms;
    \item factor bases in sieve algorithms.
\end{itemize}

This creates a general tradeoff:

\[ \boxed{\text{more memory}\quad\longleftrightarrow\quad\text{less computation}.} \]

%%%%%%%%%%%%%%%%%%%%%%%%%%%%%%%%%%%%%%%%%%%%%%%%%%%%%%%%%%%%
\subsection{Parallel Computation}
\label{subsec:number-theory-parallel-computation}
%%%%%%%%%%%%%%%%%%%%%%%%%%%%%%%%%%%%%%%%%%%%%%%%%%%%%%%%%%%%

Many number-theoretic searches can be distributed among multiple processors.

For example:

\begin{itemize}
    \item different intervals can be searched for primes;
    \item candidate factors can be tested independently;
    \item smooth-number relations can be collected in parallel;
    \item large matrix computations can be distributed.
\end{itemize}

However, not every algorithm parallelizes equally well.

Dependencies between computational steps can limit possible speedup.

%%%%%%%%%%%%%%%%%%%%%%%%%%%%%%%%%%%%%%%%%%%%%%%%%%%%%%%%%%%%
\subsection{Common Errors}
\label{subsec:computational-number-theory-common-errors}
%%%%%%%%%%%%%%%%%%%%%%%%%%%%%%%%%%%%%%%%%%%%%%%%%%%%%%%%%%%%

\subsubsection{Measuring Complexity Using the Integer Instead of Its Bit Length}

An input integer \(n\) requires only about

\[ \log_2 n \]

bits.

An algorithm polynomial in \(n\) may therefore be exponential in the actual
input length.

%%%%%%%%%%%%%%%%%%%%%%%%%%%%%%%%%%%%%%%%%%%%%%%%%%%%%%%%%%%%
\subsubsection{Computing Huge Powers Before Reducing}

To compute

\[ a^n\bmod m, \]

do not normally compute the full integer \(a^n\) first.

Reduce modulo \(m\) throughout the calculation.

%%%%%%%%%%%%%%%%%%%%%%%%%%%%%%%%%%%%%%%%%%%%%%%%%%%%%%%%%%%%
\subsubsection{Assuming a Fermat Test Proves Primality}

Composite pseudoprimes can pass Fermat tests.

Carmichael numbers can pass for every relatively prime base.

%%%%%%%%%%%%%%%%%%%%%%%%%%%%%%%%%%%%%%%%%%%%%%%%%%%%%%%%%%%%
\subsubsection{Confusing Probable Prime with Proven Prime}

A number passing randomized primality tests may be called a probable prime.

A rigorous primality proof requires a deterministic argument or verifiable
certificate.

%%%%%%%%%%%%%%%%%%%%%%%%%%%%%%%%%%%%%%%%%%%%%%%%%%%%%%%%%%%%
\subsubsection{Confusing Primality Testing with Factorization}

Knowing that an integer is composite does not automatically reveal its
factors.

%%%%%%%%%%%%%%%%%%%%%%%%%%%%%%%%%%%%%%%%%%%%%%%%%%%%%%%%%%%%
\subsubsection{Assuming Trial Division Is Efficient for Huge Integers}

Trial division is conceptually simple but scales poorly when the input has
many digits.

%%%%%%%%%%%%%%%%%%%%%%%%%%%%%%%%%%%%%%%%%%%%%%%%%%%%%%%%%%%%
\subsubsection{Using Floating-Point Arithmetic for Exact Congruences}

Congruence calculations require exact integer arithmetic.

Rounding errors can invalidate divisibility conclusions.

%%%%%%%%%%%%%%%%%%%%%%%%%%%%%%%%%%%%%%%%%%%%%%%%%%%%%%%%%%%%
\subsubsection{Assuming the Jacobi Symbol Detects Quadratic Residues}

For composite \(n\),

\[ \left(\frac an\right)=1 \]

does not guarantee that \(a\) is a square modulo \(n\).

%%%%%%%%%%%%%%%%%%%%%%%%%%%%%%%%%%%%%%%%%%%%%%%%%%%%%%%%%%%%
\subsubsection{Treating Computational Evidence as Proof}

Checking millions or billions of examples can support a conjecture but does
not automatically establish a universal theorem.

%%%%%%%%%%%%%%%%%%%%%%%%%%%%%%%%%%%%%%%%%%%%%%%%%%%%%%%%%%%%
\subsubsection{Assuming Cryptographic Hardness Is a Theorem}

Many cryptographic constructions rely on computational hardness assumptions.

The fact that no efficient classical algorithm is currently known does not
by itself prove that none exists.

%%%%%%%%%%%%%%%%%%%%%%%%%%%%%%%%%%%%%%%%%%%%%%%%%%%%%%%%%%%%
\subsection{Applications and Connections}
\label{subsec:computational-number-theory-applications}
%%%%%%%%%%%%%%%%%%%%%%%%%%%%%%%%%%%%%%%%%%%%%%%%%%%%%%%%%%%%

\begin{applicationbox}
\textbf{Cryptography.}
Modular exponentiation, prime generation, primality testing, factorization,
and discrete logarithms form the mathematical core of many cryptographic
constructions.

\medskip

\textbf{Computer Algebra.}
Exact arithmetic algorithms allow symbolic systems to manipulate enormous
integers, polynomials, algebraic numbers, and finite fields.

\medskip

\textbf{Analytic Number Theory.}
Large computations test conjectures, locate zeros of zeta and \(L\)-functions,
and investigate prime distributions.

\medskip

\textbf{Algebraic Number Theory.}
Algorithms compute rings of integers, ideal factorizations, units, class
groups, discriminants, and number-field invariants.

\medskip

\textbf{Diophantine Equations.}
Continued fractions, modular methods, lattice algorithms, and exhaustive
search help locate or exclude integer solutions.

\medskip

\textbf{Finite Fields.}
Efficient finite-field arithmetic is essential in coding theory,
cryptography, polynomial factorization, and computational algebra.

\medskip

\textbf{Algorithm Design.}
Number theory provides important examples of divide-and-reduce methods,
randomized algorithms, time--memory tradeoffs, and complexity analysis.

\medskip

\textbf{Experimental Mathematics.}
Large computations can reveal patterns and generate conjectures that later
motivate theoretical proofs.
\end{applicationbox}

%%%%%%%%%%%%%%%%%%%%%%%%%%%%%%%%%%%%%%%%%%%%%%%%%%%%%%%%%%%%
\subsection{Exercises}
\label{subsec:computational-number-theory-exercises}
%%%%%%%%%%%%%%%%%%%%%%%%%%%%%%%%%%%%%%%%%%%%%%%%%%%%%%%%%%%%

\begin{exercise}[1]
Use the Euclidean algorithm to compute

\[ \gcd(414,662). \]
\end{exercise}

\begin{exercise}[2]
Use the extended Euclidean algorithm to find integers \(x,y\) satisfying

\[ 414x+662y=\gcd(414,662). \]
\end{exercise}

\begin{exercise}[1]
Compute

\[ 19^{-1}\pmod{37}. \]
\end{exercise}

\begin{exercise}[2]
Determine whether \(24\) has a multiplicative inverse modulo \(42\).
Explain.
\end{exercise}

\begin{exercise}[1]
Use repeated squaring to compute

\[ 3^{100}\pmod7. \]
\end{exercise}

\begin{exercise}[2]
Use repeated squaring to compute

\[ 11^{123}\pmod{17}. \]
\end{exercise}

\begin{exercise}[2]
Write the exponent

\[ 157 \]

in binary and explain how its binary representation determines which powers
are multiplied in square-and-multiply.
\end{exercise}

\begin{exercise}[2]
Solve

\[ 7x\equiv5\pmod{19}. \]
\end{exercise}

\begin{exercise}[2]
Solve

\[ 12x\equiv18\pmod{30}. \]
\end{exercise}

\begin{exercise}[2]
Use the Chinese Remainder Theorem to solve

\[ x\equiv1\pmod4,\qquad x\equiv2\pmod5,\qquad x\equiv3\pmod7. \]
\end{exercise}

\begin{exercise}[1]
Use the Sieve of Eratosthenes to list all primes not exceeding \(50\).
\end{exercise}

\begin{exercise}[2]
Explain why the Sieve of Eratosthenes only needs to process primes through

\[ \sqrt N. \]
\end{exercise}

\begin{exercise}[1]
Use trial division to determine whether

\[ 137 \]

is prime.
\end{exercise}

\begin{exercise}[2]
Use Fermat's test with base \(2\) on

\[ n=15. \]
\end{exercise}

\begin{exercise}[2]
Verify that

\[ 341=11\cdot31 \]

is composite.
\end{exercise}

\begin{exercise}[3]
Verify that

\[ 2^{340}\equiv1\pmod{341}. \]

Thus show that \(341\) is a pseudoprime to base \(2\).
\end{exercise}

\begin{exercise}[2]
Factor

\[ 561 \]

and verify that it is squarefree.
\end{exercise}

\begin{exercise}[2]
Write

\[ 560=2^sd \]

with \(d\) odd in preparation for a Miller--Rabin test on \(561\).
\end{exercise}

\begin{exercise}[2]
Explain why passing one Fermat test does not establish primality.
\end{exercise}

\begin{exercise}[2]
Explain the difference between deterministic and probabilistic primality
testing.
\end{exercise}

\begin{exercise}[2]
If \(10\) independent Miller--Rabin rounds are performed, use the standard
bound to estimate the maximum false-acceptance probability.
\end{exercise}

\begin{exercise}[1]
Factor

\[ 8051 \]

by trial division.
\end{exercise}

\begin{exercise}[2]
Use Fermat factorization to factor

\[ 2021. \]
\end{exercise}

\begin{exercise}[2]
Explain why Fermat factorization is especially effective when the two factors
of \(n\) are close together.
\end{exercise}

\begin{exercise}[2]
For

\[ f(x)=x^2+1\pmod{91}, \]

compute the first several terms beginning with \(x_0=2\).
\end{exercise}

\begin{exercise}[3]
Explain how a collision modulo an unknown factor can cause

\[ \gcd(x_i-x_j,n) \]

to reveal that factor in Pollard's rho method.
\end{exercise}

\begin{exercise}[1]
Determine whether each number is \(5\)-smooth:

\[ 72,\qquad90,\qquad98,\qquad125,\qquad150. \]
\end{exercise}

\begin{exercise}[2]
Write the parity exponent vectors over \(\mathbf F_2\) for

\[ 12=2^2\cdot3,\qquad18=2\cdot3^2,\qquad24=2^3\cdot3. \]
\end{exercise}

\begin{exercise}[3]
Explain why finding a linear dependence among exponent vectors modulo \(2\)
can produce a perfect-square product.
\end{exercise}

\begin{exercise}[1]
Given

\[ n=2^4\cdot3^2\cdot5, \]

compute

\[ \tau(n),\qquad\sigma(n),\qquad\varphi(n),\qquad\mu(n). \]
\end{exercise}

\begin{exercise}[2]
Use Euler's criterion to determine whether \(5\) is a quadratic residue
modulo \(11\).
\end{exercise}

\begin{exercise}[2]
Compute

\[ \left(\frac{7}{19}\right) \]

using quadratic reciprocity.
\end{exercise}

\begin{exercise}[2]
Compute the Jacobi symbol

\[ \left(\frac{5}{21}\right). \]
\end{exercise}

\begin{exercise}[2]
Explain why a Jacobi symbol equal to \(1\) does not necessarily imply a
quadratic residue when the denominator is composite.
\end{exercise}

\begin{exercise}[2]
Find a square root of \(4\) modulo \(11\) using

\[ x\equiv a^{(p+1)/4}\pmod p. \]
\end{exercise}

\begin{exercise}[1]
Solve the discrete logarithm

\[ 2^x\equiv8\pmod{11}. \]
\end{exercise}

\begin{exercise}[2]
Solve

\[ 3^x\equiv13\pmod{17}. \]
\end{exercise}

\begin{exercise}[3]
Explain why baby-step giant-step reduces a search of size \(N\) to roughly

\[ \sqrt N \]

baby and giant steps.
\end{exercise}

\begin{exercise}[2]
Compute the continued fraction of

\[ \frac{43}{17}. \]
\end{exercise}

\begin{exercise}[2]
Find the first four convergents of

\[ [1;2,2,2,\ldots]. \]
\end{exercise}

\begin{exercise}[3]
Explain why continued fractions are useful for Pell equations.
\end{exercise}

\begin{exercise}[2]
Verify that

\[ 3^2-2(2^2)=1. \]
\end{exercise}

\begin{exercise}[2]
Suppose

\[ p=11,\qquad q=17. \]

Compute

\[ n=pq \]

and

\[ \varphi(n). \]
\end{exercise}

\begin{exercise}[2]
Using the previous exercise, choose

\[ e=7 \]

and compute an inverse

\[ d\equiv e^{-1}\pmod{\varphi(n)}. \]
\end{exercise}

\begin{exercise}[3]
Explain why knowing \(p\) and \(q\) makes computing \(\varphi(pq)\) easy.
\end{exercise}

\begin{exercise}[3]
Explain why efficient modular exponentiation is essential for RSA-style
arithmetic.
\end{exercise}

\begin{exercise}[2]
Explain the distinction between a computationally difficult problem and a
mathematically impossible problem.
\end{exercise}

\begin{exercise}[2]
Give one example of a randomized number-theoretic algorithm and explain what
role randomness plays.
\end{exercise}

\begin{exercise}[2]
Classify the general behavior of probabilistic Miller--Rabin testing as more
closely related to a Monte Carlo or Las Vegas algorithm.
\end{exercise}

\begin{exercise}[2]
Explain why verifying

\[ n=ab \]

is usually much easier than discovering the factors \(a\) and \(b\).
\end{exercise}

\begin{exercise}[3]
Suppose a conjecture has been verified computationally for every

\[ n\leq10^{15}. \]

Explain why this does not automatically prove the conjecture for every
positive integer.
\end{exercise}

\begin{exercise}[3]
Explain the computational chain

\[
\text{Euclidean algorithm}
\longrightarrow
\text{modular inverse}
\longrightarrow
\text{modular arithmetic}
\longrightarrow
\text{cryptographic computation}.
\]
\end{exercise}

\begin{exercise}[3]
Explain the computational relationship

\[
\text{prime generation}
\longrightarrow
\text{primality testing}
\longrightarrow
\text{large integer arithmetic}
\longrightarrow
\text{public-key mathematics}.
\]
\end{exercise}

%%%%%%%%%%%%%%%%%%%%%%%%%%%%%%%%%%%%%%%%%%%%%%%%%%%%%%%%%%%%
\subsection{Conceptual Summary}
\label{subsec:computational-number-theory-summary}
%%%%%%%%%%%%%%%%%%%%%%%%%%%%%%%%%%%%%%%%%%%%%%%%%%%%%%%%%%%%

Computational number theory studies how arithmetic problems can be solved
algorithmically.

The Euclidean algorithm efficiently computes

\[ \gcd(a,b), \]

while the extended Euclidean algorithm computes Bézout coefficients and
modular inverses.

Repeated squaring allows powers

\[ a^n\bmod m \]

to be evaluated using only logarithmically many exponentiation stages.

The Chinese Remainder Theorem provides an algorithmic method for combining
compatible congruence information.

Prime-generation algorithms such as the Sieve of Eratosthenes efficiently
produce collections of primes.

Primality testing asks whether an integer is prime, while factorization asks
for its prime factors.

Fermat's test motivates probabilistic primality methods but is defeated by
pseudoprimes and Carmichael numbers.

Miller--Rabin provides a much stronger practical test.

Integer factorization progresses from trial division and Fermat's method to
Pollard's rho, quadratic-sieve methods, and the number field sieve.

Quadratic residues can be studied computationally through Euler's criterion,
quadratic reciprocity, Legendre symbols, Jacobi symbols, and modular
square-root algorithms.

Discrete logarithms reverse exponentiation in finite cyclic groups and lead
to algorithms such as baby-step giant-step and Pollard's rho.

Continued fractions connect the Euclidean algorithm with rational
approximation and Pell equations.

Modern cryptographic mathematics depends heavily on efficient prime
generation, modular inversion, modular exponentiation, factorization
problems, and discrete logarithms.

Computational experiments can discover patterns and test enormous numbers of
cases, but computation and proof play different mathematical roles.

The central principle is:

\[
\boxed{
\text{knowing that an arithmetic object exists is different from knowing how
to compute it efficiently}.
}
\]

%%%%%%%%%%%%%%%%%%%%%%%%%%%%%%%%%%%%%%%%%%%%%%%%%%%%%%%%%%%%
\subsection{Section Connections}
\label{subsec:computational-number-theory-connections}
%%%%%%%%%%%%%%%%%%%%%%%%%%%%%%%%%%%%%%%%%%%%%%%%%%%%%%%%%%%%

\chapterconnection{
Computational Number Theory turns the theoretical structures of Elementary
Number Theory into explicit algorithms. The Euclidean algorithm, modular
inverses, fast exponentiation, congruence solving, primality testing, and
factorization provide the basic computational toolkit. Analytic Number
Theory motivates large-scale computations involving primes and arithmetic
functions, while Algebraic Number Theory leads to algorithms for number
fields, ideals, class groups, and units. Linear Algebra enters modern
factorization through exponent vectors and sparse systems over finite
fields. Probability motivates randomized primality and factorization
algorithms. Continued fractions connect computation with Diophantine
Approximation. Discrete logarithms and integer factorization connect number
theory with cryptography. More general questions of numerical efficiency,
bit complexity, symbolic computation, and high-performance algorithms are
developed further in Numerical and Computational Mathematics.
}

%%%%%%%%%%%%%%%%%%%%%%%%%%%%%%%%%%%%%%%%%%%%%%%%%%%%%%%%%%%%
% SECTION SOURCE TRACKING
%%%%%%%%%%%%%%%%%%%%%%%%%%%%%%%%%%%%%%%%%%%%%%%%%%%%%%%%%%%%

%------------------------------------------------------------
% 2.4 Computational Number Theory
%------------------------------------------------------------
%
% Source basis:
%   - Standard undergraduate and introductory computational
%     number theory and algorithmic number theory.
%
% Used for:
%   - algorithmic viewpoint and input size;
%   - bit length and basic complexity;
%   - Euclidean and extended Euclidean algorithms;
%   - modular inverses;
%   - repeated squaring and square-and-multiply;
%   - computational Chinese Remainder Theorem;
%   - Sieve of Eratosthenes;
%   - primality testing;
%   - Fermat pseudoprimes and Carmichael numbers;
%   - Miller--Rabin testing;
%   - deterministic primality testing;
%   - primality certificates;
%   - integer factorization;
%   - Fermat factorization;
%   - Pollard's rho;
%   - smooth numbers and quadratic-sieve ideas;
%   - number field sieve overview;
%   - arithmetic-function computation;
%   - Legendre and Jacobi symbol computation;
%   - modular square roots;
%   - discrete logarithms;
%   - baby-step giant-step;
%   - continued fractions;
%   - Pell equations;
%   - RSA and Diffie--Hellman mathematical foundations;
%   - randomized algorithms;
%   - computational experiments and verification;
%   - exact and large-integer arithmetic;
%   - time--memory and parallel-computation considerations.
%
% Notes:
%   - Divisibility, congruences, Euler's theorem, quadratic
%     residues, and the Chinese Remainder Theorem are taught
%     canonically in Elementary Number Theory.
%   - General asymptotic complexity belongs more broadly to
%     discrete mathematics and computer science.
%   - Number-field algorithms are developed after the
%     theoretical structures of Algebraic Number Theory.
%   - Advanced Diophantine algorithms are deferred to
%     Diophantine Approximation and Diophantine Geometry.
%   - Cryptographic constructions are included only to
%     illustrate the mathematical role of the algorithms;
%     secure implementation is outside the scope of this text.
%
%%%%%%%%%%%%%%%%%%%%%%%%%%%%%%%%%%%%%%%%%%%%%%%%%%%%%%%%%%%%